\section*{Supplementary Material}
\addcontentsline{toc}{section}{Supplementary Material}
\renewcommand{\thesection}{S\arabic{section}}
\setcounter{section}{0}

\section{Full EMOD d6 SOTA recipe}
\label{sup:emod_recipe}

Our FACED SOTA uses the following recipe (inherited from prior mechanism work):
\begin{itemize}[leftmargin=*]
\item \textbf{Backbone.} EMOD axial-transformer (\cite{chen2026emod}).
\item \textbf{Depth extension.} Pretrained EMOD has 3 axial layers; we
increase \texttt{transformer\_depth=6} and load the pretrained checkpoint
with \texttt{strict=False}, keeping layers 0--2 pretrained and randomly
initializing layers 3--5 (``progressive stacking'' of
\cite{gong2024progressive}).
\item \textbf{Augmentation.} CBraMod-style 4-op augmentation
(jitter / noise / channel-dropout / amplitude-scale) with $p_\mathrm{aug}=0.6$.
\item \textbf{Label smoothing.} $\mathrm{LS}=0.1$ on the cross-entropy loss.
\item \textbf{Knowledge distillation.} Auxiliary KD term with $\lambda=0.1$,
$\tau=0.1$ and a fixed random orthonormal 9$\times$1024 prototype tensor
(we showed earlier that replacing this with LLM-derived vectors does not
change the result within noise).
\item \textbf{Optimizer.} Adam with $\mathrm{lr}=10^{-3}$, weight decay
$10^{-4}$, backbone-LR multiplier $0.5$.
\item \textbf{Schedule.} Cosine annealing with $T_\mathrm{max}=100$ (group A)
or $T_\mathrm{max}=150$ (group B), $\mathrm{lr}_{\mathrm{min}}=10^{-6}$.
\item \textbf{Batch size.} 64.
\item \textbf{Seeds.} $\{2025, 42, 123, 456, 789\}$ for each group.
\end{itemize}

\section{Full per-seed results}
\label{sup:per_seed}

\begin{table}[h]
\centering
\caption{Per-seed test BAcc for FACED d6 100ep and 150ep trainings.}
\label{tab:perseed_faced}
\begin{tabular}{lcc}
\toprule
Seed & 100ep BAcc & 150ep BAcc \\
\midrule
2025 & 0.6638 & 0.6565 \\
42   & 0.6549 & 0.6526 \\
123  & 0.6479 & 0.6573 \\
456  & 0.6625 & 0.6486 \\
789  & 0.6612 & 0.6755 \\
\midrule
mean & 0.6581 & 0.6581 \\
\bottomrule
\end{tabular}
\end{table}

\section{Alternative endpoints}
\label{sup:endpoints}

We ran small ablations varying $T_B/T_A$. Results on CIFAR-10:

\begin{table}[h]
\centering
\caption{Effect of $T_B$ on CIFAR-10 mixed-length ensemble gain (with
$T_A=50$ fixed).}
\label{tab:endpoint_ablation}
\begin{tabular}{lccc}
\toprule
$T_B$ & 5-seed $T_A$ ens. & Mixed ensemble & Gain over $T_A$ ens. \\
\midrule
50 (same, control) & 0.9209 & 0.9209 & 0.000 \\
65  & TBD & TBD & TBD \\
75 (main)  & 0.9209 & 0.9253 & +0.0044 \\
100 & TBD & TBD & TBD \\
\bottomrule
\end{tabular}
\end{table}

\section{Training logs for 100ep vs 150ep FACED d6}
\label{sup:training_logs}

Full per-epoch val BACC curves are available in the supplementary data
release (\texttt{training\_log\_analysis.json}). The key observation is
that the 100ep and 150ep trajectories reach nearly identical val BACC by
their respective final epochs ($0.6771$ vs.\ $0.6776$), but the 150ep
trajectory is behind at epoch 99 ($0.6722$) because the cosine schedule
still has $\sim\!25\%$ of the base LR remaining.

\section{Error-correlation analysis}
\label{sup:correlation}

\begin{table}[h]
\centering
\caption{Error correlation (Cohen $\kappa$ on error indicators) within and
between the 100ep and 150ep groups. Lower $\kappa$ means more independent
errors, which drives ensemble gain.}
\label{tab:error_kappa}
\begin{tabular}{lc}
\toprule
Comparison & Cohen $\kappa$ \\
\midrule
Within 100ep group (10 pairs) & 0.6935 \\
Within 150ep group (10 pairs) & 0.6786 \\
Cross-group 100ep $\times$ 150ep (25 pairs) & 0.7024 \\
\bottomrule
\end{tabular}
\end{table}
